\begin{statementbox}
	\textbf{\textcolor{CStatement}{条件}}
	\begin{description}[style=nextline, leftmargin=2.8em, labelsep=0.5em, itemsep=0.35em]
		\item[\textbf{\textcolor{CStatement}{条件1（圆锥曲线非退化）：}}]
			给定非退化的圆锥曲线 $\Gamma$（标准椭圆/双曲线/抛物线，或一般二次曲线 $Ax^2+Bxy+Cy^2+Dx+Ey+F=0$）。
		\item[\textbf{\textcolor{CStatement}{条件2（给定一点P）：}}]
			平面上一点 $P(x_0,y_0)$，可在曲线外、上或内。
	\end{description}
	\textbf{\textcolor{CStatement}{结论}}
	\begin{description}[style=nextline, leftmargin=2.8em, labelsep=0.5em, itemsep=0.45em]
		\item[\textbf{\textcolor{CStatement}{结论一（标准椭圆双曲线）：}}]
			\textbf{\textcolor{CStatement}{直观描述：}}
			将点的坐标代入曲线方程中的平方项一次化。
			\[
				\frac{x_0x}{a^2}\pm\frac{y_0y}{b^2}=1.
			\]
			其中椭圆取 $+$，双曲线取 $-$。
		\item[\textbf{\textcolor{CStatement}{结论二（抛物线情形）：}}]
			\textbf{\textcolor{CStatement}{直观描述：}}
			将点坐标代入抛物线方程中的变量一次化。
			\[
				y_0y=p(x+x_0).
			\]
		\item[\textbf{\textcolor{CStatement}{结论三（一般二次曲线）：}}]
			\textbf{\textcolor{CStatement}{直观描述：}}
			对二次项和一次项分别做对称替换。
			\[
				A x_0 x + \frac{B}{2}(x_0 y + y_0 x) + C y_0 y + \frac{D}{2}(x+x_0) + \frac{E}{2}(y+y_0) + F = 0.
			\]
	\end{description}
	\textbf{\textcolor{CStatement}{推广结论（切点弦应用）}}
	\textbf{\textcolor{CStatement}{直观描述：}}
	当从 $P$ 存在两条具有不同有限实切点的切线时，极线恰为切点弦所在直线（即从 $P$ 所作两条切线的切点连线）。

\medskip
\textbf{有限平面中的定义域。} 有心曲线的中心须排除；在标准式中 $P=(0,0)$ 给出 $0=1$，
不是有限直线。一般式要求一次项系数不同时为零。
极线可没有实交点；只有确有两不同有限实切点时，才能把它称作实切点弦。
\end{statementbox}
